\documentclass{apcalc}
\title{Unit 8: Applications of Integration}
\author{Paul Wang}
\date{November 2026}
\volume{Study Guide 8 of 10 (topics 8.1--8.13)}
\shorttitle{Unit 8 --- Study Guide}
\begin{document}
\maketitle

\begin{abstract}
Unit 8 turns the definite integral into averages, distances, amounts, areas, volumes and (in BC) arc lengths. It carries 5--10\% of the multiple-choice section and supplies at least one free-response question every year: an area-and-volume question on a region between two curves, and a rate question (``water enters at $E(t)$ and leaves at $L(t)$'') in the calculator part. The skill that earns the points is setting up the integral correctly: limits, integrand and, for volumes, the radius measured from the axis. This guide states the unit in eleven rules (R8.1--R8.11) across five sections. The \emph{Unit 8 Practice Set} uses the same section numbers.
\end{abstract}
\keywords{average value, displacement, total distance, net change, area between curves, cross sections, disc method, washer method, axis of revolution, arc length.}

\section*{Router}
\begin{router}
``average value of $f$ on $[a,b]$'' & $\frac{1}{b-a}\int_a^bf$ (not the average rate of change) & R8.1 \\
``total distance traveled'' & $\int_a^b|v(t)|\,dt$ & R8.2 \\
``position at $t=b$'' & $x(a)+\int_a^bv(t)\,dt$ & R8.2 \\
rate in and rate out & amount $=A_0+\int(\text{in}-\text{out})$; max where in $=$ out changes sign & R8.3 \\
area between $y=f$ and $y=g$ & $\int(\text{top}-\text{bottom})\,dx$ or $\int(\text{right}-\text{left})\,dy$ & R8.4--R8.6 \\
known cross sections & $\int A(x)\,dx$ with the side $s=$ top $-$ bottom & R8.7 \\
revolve, region touches the axis & disc: $\pi\int R^2$ & R8.8 \\
revolve, gap between region and axis & washer: $\pi\int(R^2-r^2)$ & R8.9 \\
axis $y=k$ or $x=k$ & radius $=$ distance from the axis: $|k-y|$ or $|k-x|$ & R8.10 \\
``length of the curve'' & $\int_a^b\sqrt{1+(f')^2}\,dx$ & R8.11 \\
\end{router}

\section{Average value (8.1)}

\begin{rulebox}{R8.1}{Average value of a function}
$f_{\text{avg}}=\dfrac{1}{b-a}\displaystyle\int_a^bf(x)\,dx$. By the Mean Value Theorem for integrals, a continuous $f$ equals its average value at least once on $[a,b]$. Do not confuse with the average \emph{rate of change} $\frac{f(b)-f(a)}{b-a}$; the average value of $f'$ equals the average rate of change of $f$.
\end{rulebox}

\section{Motion and accumulation in context (8.2--8.3)}

\begin{rulebox}{R8.2}{Displacement, distance and position}
Displacement on $[a,b]$: $\int_a^bv(t)\,dt$. Total distance: $\int_a^b|v(t)|\,dt$ (on the calculator directly, or split at the zeros of $v$ by hand). Position: $x(b)=x(a)+\int_a^bv(t)\,dt$. Similarly $v(b)=v(a)+\int_a^ba(t)\,dt$.
\end{rulebox}

\begin{example}{distance versus displacement}
$v(t)=t^2-4t+3=(t-1)(t-3)$ on $[0,4]$. Displacement $\int_0^4v=\frac{64}{3}-32+12=\frac43$. Distance: $\int_0^1v-\int_1^3v+\int_3^4v=\frac43+\frac43+\frac43=4$. \answer{displacement $\frac43$, distance 4.}
\end{example}

\begin{rulebox}{R8.3}{Rates in, rates out}
Amount at time $T$: $A(T)=A(0)+\int_0^T\big(\text{in}(t)-\text{out}(t)\big)\,dt$. The amount is increasing when in $>$ out. For an absolute maximum on $[0,T]$, find where $\text{in}-\text{out}$ changes sign, then compare the amounts at those times and at the endpoints (candidates test). Units: rate units times time units.
\end{rulebox}

\begin{frqtip}
In Part A these questions are calculator active: store the functions, write the integral expression with limits, and give answers to three decimal places. ``$\int_0^{6}E(t)\,dt\approx 52.874$ gallons'' earns the point only if the integral is written down.
\end{frqtip}

\section{Area between curves (8.4--8.6)}

\begin{rulebox}{R8.4}{Area with respect to $x$}
$\displaystyle A=\int_a^b\big(\text{top}(x)-\text{bottom}(x)\big)\,dx$, with $a$, $b$ the $x$-values where the curves meet (or the given bounds).
\end{rulebox}

\begin{rulebox}{R8.5}{Area with respect to $y$}
When the region is bounded on the left and right by functions of $y$ (for example $x=y^2$), use $\displaystyle A=\int_c^d\big(\text{right}(y)-\text{left}(y)\big)\,dy$.
\end{rulebox}

\begin{rulebox}{R8.6}{Curves that cross}
If the curves cross inside $[a,b]$, split at each intersection so that each piece is top minus bottom, or integrate $|f-g|$ on the calculator.
\end{rulebox}

\begin{example}{area between a line and a parabola}
\begin{center}
\begin{tikzpicture}
\begin{axis}[width=6cm,height=4.6cm,axis lines=middle,xmin=-0.5,xmax=3.5,ymin=-0.5,ymax=4.8,xtick={1,2,3},ytick={1,2,3,4},xlabel={$x$},ylabel={$y$},
 every axis x label/.style={at={(ticklabel* cs:1)},anchor=west},every axis y label/.style={at={(ticklabel* cs:1)},anchor=south}]
\addplot[name path=A,thick,domain=0:3]{2*x-0*x} node[pos=0.9,above left]{$y=2x$};
\addplot[name path=B,thick,domain=0:3,samples=60]{x^2} ;
\addplot[gray!35] fill between[of=A and B,soft clip={domain=0:2}];
\node at (axis cs:2.9,4.2) {$y=x^2$};
\end{axis}
\end{tikzpicture}
\end{center}
Region $R$ between $y=2x$ and $y=x^2$: they meet at $x=0$ and $x=2$, the line on top. $A=\int_0^2(2x-x^2)\,dx=4-\frac83=\frac43$. \answer{$\frac43$.}
\end{example}

\section{Volumes (8.7--8.12)}

\begin{rulebox}{R8.7}{Known cross sections}
$V=\int_a^bA(x)\,dx$ where $A(x)$ is the area of the cross section perpendicular to the $x$-axis. With side (or diameter) $s(x)=\text{top}-\text{bottom}$: square $s^2$; semicircle $\frac\pi8s^2$; equilateral triangle $\frac{\sqrt3}{4}s^2$; isosceles right triangle with leg $s$: $\frac12s^2$; rectangle of height $h\cdot s$: $hs^2$.
\end{rulebox}

\begin{rulebox}{R8.8}{Disc method}
Revolving a region that touches the axis: $V=\pi\int_a^bR(x)^2\,dx$, with $R$ the distance from the axis to the curve.
\end{rulebox}

\begin{rulebox}{R8.9}{Washer method}
With a gap between the region and the axis: $V=\pi\int_a^b\big(R(x)^2-r(x)^2\big)\,dx$, outer radius $R$, inner radius $r$. Square each radius separately: $\pi\int(R-r)^2$ is wrong.
\end{rulebox}

\begin{rulebox}{R8.10}{Other axes}
Measure radii from the axis. About $y=-1$: radius to the curve $y=f(x)$ is $f(x)-(-1)=f(x)+1$. About $y=k$ above the region: $k-f(x)$. About a vertical line $x=k$, integrate in $y$ with radii $|k-x(y)|$.
\end{rulebox}

\begin{example}{washer about a horizontal line}
Revolve the region $R$ above about $y=-1$: outer radius $2x+1$, inner radius $x^2+1$.
$V=\pi\int_0^2\big((2x+1)^2-(x^2+1)^2\big)\,dx=\pi\int_0^2(-x^4+2x^2+4x)\,dx=\pi\left(-\frac{32}{5}+\frac{16}{3}+8\right)=\frac{104\pi}{15}$. \answer{$\frac{104\pi}{15}$.}
\end{example}

\section{Arc length (8.13)}

\begin{rulebox}{R8.11}{Arc length (BC)}
For $y=f(x)$ on $[a,b]$ with $f'$ continuous: $L=\displaystyle\int_a^b\sqrt{1+\big(f'(x)\big)^2}\,dx$. For $x=g(y)$: $\int_c^d\sqrt{1+g'(y)^2}\,dy$. Most arc lengths need a calculator; set up the integral first. (Parametric arc length is in Unit 9.)
\end{rulebox}

\begin{example}{an exact arc length}
$y=\frac23x^{3/2}$ on $[0,3]$: $y'=x^{1/2}$, $L=\int_0^3\sqrt{1+x}\,dx=\frac23\left[(1+x)^{3/2}\right]_0^3=\frac23(8-1)=\frac{14}{3}$.
\end{example}

\section{Common errors}
\begin{itemize}
\item Computing $\frac{f(b)-f(a)}{b-a}$ for an average value (R8.1).
\item Reporting displacement when asked for distance (R8.2).
\item Forgetting the initial amount, or using in $+$ out (R8.3).
\item Bottom minus top, or not splitting at crossings (R8.4--R8.6).
\item $\pi\int(R-r)^2$ instead of $\pi\int(R^2-r^2)$; radius measured from the $x$-axis when the axis is $y=k$ (R8.9--R8.10).
\item Forgetting to square $f'$ in arc length (R8.11).
\end{itemize}

\end{document}
